% Generated draft: requires mathematical and source review.
\documentclass[11pt]{article}
\usepackage[margin=25mm]{geometry}
\usepackage{fontspec}
\setmainfont{Latin Modern Roman}
\usepackage{amsmath,amssymb,mathtools,unicode-math}
\setmathfont{Latin Modern Math}
\usepackage{enumitem,needspace,xurl,hyperref}
\setlength{\parindent}{0pt}
\setlength{\parskip}{5pt}
\setlength{\emergencystretch}{3em}
\newcommand{\N}{\mathbb{N}}
\newcommand{\Z}{\mathbb{Z}}
\newcommand{\Q}{\mathbb{Q}}
\newcommand{\R}{\mathbb{R}}
\newcommand{\C}{\mathbb{C}}
\newcommand{\F}{\mathbb{F}}
\newcommand{\A}{\mathbb{A}}
\newcommand{\PP}{\mathbb{P}}
\newcommand{\E}{\mathbb{E}}
\newcommand{\eps}{\varepsilon}
\newcommand{\dd}{\,\mathrm{d}}
\begin{document}
\section*{UnsolvedMath ID: 20002006}
\textbf{Source problem code:} catalogue entry
\subsection{Definitions and mathematical statement}
The original Section C.1, ``Relevant lines of research,'' is a thematic list rather than a quantified proposition. It lists: basic definitions; computational issues; higher-codimension amoebas; example families; abstract data types in tropical and idempotent calculus; recognition problems; applications to complex algebraic geometry, real algebraic geometry, dynamical systems, differential equations, optimization and control, representation theory, number theory and statistical mechanics; and tropical representation theory. No domain, desired output or theorem is attached to these headings.

The aid supplies a separate precise recognition problem inspired by that list. In \(\R^2\), take distinct rays with primitive integer generators \(u_1,\ldots,u_m\in\Z^2\setminus\{0\}\), in cyclic order, and independently specified positive integer tropical cell weights \(w_1,\ldots,w_m\). Suppose the directions span \(\R^2\). The reported criterion is
\[
 \text{the weighted rays form the codimension-one skeleton of the normal fan
 of a lattice polygon}
 \quad\Longleftrightarrow\quad
 \sum_{i=1}^m w_i u_i=0.
\]
Here the normal fan uses outward normals and the max-plus convention; the weight of a ray is the lattice length of its dual polygon edge. Equivalently this is a constant-coefficient one-vertex planar tropical hypersurface.

Let \(J(x,y)=(-y,x)\). For accepted full-rank data the explicit polygon is reconstructed by
\[
 p_0=0,\qquad p_k=\sum_{i=1}^k w_iJu_i\quad(1\leq k\leq m),
 \qquad p_m=p_0.
\]
Its edge directions, lengths and convex cyclic order give the unique Newton polygon up to translation. A corresponding tropical polynomial can use its vertices with equal coefficients. The finite certificate checks integer primitivity, distinct rays, cyclic order, rank two, positive integer weights, and the two-coordinate balance sum, followed by these exact prefix sums.

For rank-one nonempty balanced distinct primitive-ray data, there are exactly two antipodal rays \(u,-u\), with the same weight \(w>0\). They form the multiplicity-\(w\) binomial line, dual to the lattice segment \([0,wJu]\). This degenerate case is not represented as a full-dimensional polygon.

Normalization conventions matter. If input uses a nonprimitive representative \(v_i=d_i u_i\), \(d_i>0\), and \(w_i\) independently denotes the tropical cell weight, replace \(v_i\) by \(u_i\) while retaining \(w_i\). The different rule \(w_i\mapsto d_iw_i\) is appropriate only when the raw product \(w_iv_i\) was explicitly intended as the balancing vector. Confusing these encodings changes the recognized tropical cycle. The criterion is a reported special synthesis; it is not a theorem that every source research direction has been solved, nor recognition for arbitrary nonconstant-coefficient tropical varieties.
\subsection{Short English statement}
The source is a list of research directions. Its aid develops one recognition example: primitive weighted planar rays give a lattice-polygon normal fan exactly when they balance, with an explicit polygon reconstruction. Rank-one data give a weighted binomial line, and ray normalization must preserve the declared cell-weight convention.
\subsection{Sources}
\begin{itemize}
\item Local UnsolvedMath registry, permanent ID \(20002006\), source code \texttt{catalogue entry}; canonical collection \texttt{aim-geometry-notes.json}, notes[343], zero-based source index \(343\), tag ``section.''
\item AIM workshop ``Amoebas and tropical geometry,'' 23--26 October 2003, Section C.1, ``Relevant lines of research'': \url{https://aimath.org/WWN/amoebas/amoebas.pdf}. No PDF was opened in this drafting session.
\item Supplied research attempt \(1\), result \texttt{context\_only}, run status \texttt{not\_a\_problem}, rigor label ``proved,'' novelty confidence ``low.'' Registry status ``open,'' classification \texttt{OPEN-TRIAGE}, suggested difficulty \texttt{L2}, updated 14 August 2026. The recognition theorem is the aid's mathematical synthesis, not an original question stated in C.1.
\item Original remarks and literature fields are empty; no proposer, proposed year or authoritative external citation for the recognition criterion is supplied.
\end{itemize}
\end{document}
